% Copyright 2026 Yilin Liu and GPT-6 Astra Pro. Apache-2.0.
% Curated connective exposition of the checkpoint's rectangular extraction,
% direct mixed-crossing transfer, and integral full-curve numerical consequence.
\section{Integral rectangular extraction and shared evaluation}
\label{sec:rectangular-hurwitz}
For integer dimensions $p,q\ge2$, choose the least $h,E$ with
$p(h)\ge p$ and $p(E)\ge q$, where $p(t)$ denotes the partition
function in this paragraph only. Use any $p$ partitions of $h$ and $q$
partitions of $E$ in Theorem~\ref{thm:integral-crossing}. Restriction to
these coordinates preserves its integer crossing identity. For each
piece select its smaller rank factorization. The number of resulting
rank-one terms is at most
\[
 (2(E+h)+1)(2h+1)p(E+h).
\]
The partition estimate
$\log p(t)=\pi\sqrt{2t/3}+O(\log(t+2))$
\cite{hardyramanujan1918} therefore gives
\[
 P\le\exp\!\left(\sqrt{(\log p)^2+(\log q)^2}\right)
      (\log(pq+2))^{O(1)}.
\]
Here logarithms are natural. The rank-one terms have the two orientations
of \eqref{eq:direct-mixed-crossing}; the selected factorizations are
integer, with normalization one. No conversion to squares or division
by two is made. Their entries have polynomial logarithmic bit length
and are explicitly computable by
Lemma~\ref{lem:algebraic-coefficient-computation}.

Fix $a,b>0$ and $L\ge\max\{a,b\}$, and set
$s=\sqrt{a^2+b^2}$. Choose base dimensions
$p=n^{a+o(1)}$, $q=n^{b+o(1)}$. Then
$P=n^{s+o(1)}$. Partition each target variable list into
$U=n^{L-a+o(1)}$ and $V=n^{L-b+o(1)}$ blocks, respectively.
Replicate the same integer crossing decomposition for every block pair,
using the same pair at both repeated occurrences. This yields an exact
mixed crossing for $ID_{n^{L+o(1)}}$; padding adds zero coordinates.

For a factor family with coefficient matrices $A_j\in\mathbb Z^{p\times q}$,
all forms
\[
 (X^{(r)})^{\mathsf T}A_jY^{(t)},
 \quad j\in[P],\ r\in[U],\ t\in[V],
\]
are computed by two shared matrix products. Stack the $A_j$ into an
$(Pp)\times q$ constant matrix and multiply by the $q\times V$ matrix
of $Y$ blocks. Reshape the result to $p\times(PV)$, then multiply on
the left by the $U\times p$ matrix of $X$ blocks. Both factor families
require only a constant number of such calls. If their integer matrix
schemes have full arithmetic exponents $W_1,W_2$, the two shapes are
\[
 (s+a,b,L-b),\qquad (L-a,a,s+L-b),
\]
and Lemma~\ref{lem:direct-mixed-crossing} gives exponent
$\max\{W_1,W_2\}/L$ for the ordered identity. All $PUV$ output forms
are charged by the matrix schemes' output dimension, so insertion's
linear overhead is absorbed. Matrix products are used first as
commutative bilinear circuits in formal variables; ordered insertion
then preserves the required word order. The integer scheme, including
all scalar-linear operations, is fixed before evaluating any target
associative ring.

When $a+b=L$, the shapes are permutations of
$(s+a,b,a)$. Tensor-leg permutations of an integer bilinear matrix
identity give valid ordered schemes: interpret the first chosen leg
as the left input and the second as the right input, without commuting
target-ring values. This is the specialization used below.

\section{The selected constant-free ordered circuit bound}
\begin{theorem}[One integer circuit family with exponent $1.719516$]
\label{thm:selected-crossing}
Using the integral source constructions and weighted-profile theorem
specified in Theorem~\ref{thm:integral-full-rectangular}, there is one
family of noncommutative arithmetic circuits over $\mathbb Z$ computing
\[
 ID_N=\sum_{i,j=1}^{N}X_iY_jX_iY_j
\]
with $O(N^{1.719516})$ gates and only scalar constants $0,1,-1$.
The same chosen circuit evaluates correctly in every associative
unital ring, including characteristic two and rings with zero divisors.
The bound counts arithmetic gates; it does not bound the bit cost of
arbitrary ring-element evaluations or give a circuit-generation time
matching this gate exponent.
\end{theorem}
\begin{proof}
Set
\[
 a=\frac{167}{400},\qquad b=\frac{233}{400},\qquad L=1,
 \qquad s=\frac{\sqrt{82178}}{400}.
\]
Both shared matrix shapes are permutations of $(s+a,b,a)$. Put
\[
 R=\frac{286667054264}{10^9},\qquad A=a+R/400.
\]
Integer squaring gives $R^2>82178$. Padding the largest dimension
therefore permits the rational shape $(A,b,a)$, with $A>b>a>0$.
It lies in the curved region $a(A+b)>Ab$, and exact rational
substitution gives
\[
 \frac{1719516}{10^6}
 -\frac{(Ab+ba+aA)^2}{4Aba}
 =\frac{358038466018273}{1379112402223945625000}>0.
\]
Theorem~\ref{thm:integral-full-rectangular} supplies one integer
ordered matrix scheme at this rational shape, with any fixed
positive exponent slack. Reserve part of the strict displayed gap
for that slack. The shared-products construction and
Lemma~\ref{lem:direct-mixed-crossing} supply the ordered $ID_N$
circuit with the same exponent, apart from the partition factors,
rounding and padding. Their polynomial-logarithmic overhead fits
another part of the same gap.

There are two kinds of scalar constants. Crossing entries use
$h,E=O((\log(N+2))^2)$ and have bit length
$(\log(N+2))^{O(1)}$ by
Lemma~\ref{lem:algebraic-coefficient-computation}. For the matrix
scheme, fix the finite integer witness after fixing the positive
slack. Its finite collection of integer constants is independent of
$N$. Recursive evaluation, bihomogenization, scalar-linear
transposition and insertion use only these constants and $0,1,-1$.
If a scalar-only subcircuit appears, retain it instead of replacing
it by a new unanalyzed scalar leaf.

An integer of bit length $d$ is built by $O(d+1)$ doublings, additions
and a possible negation from $0,1,-1$. Replacing every scalar leaf
therefore costs only a polynomial-logarithmic factor. Reserve the
remaining strict gap for this cost. All constructions are identities
in the free associative algebra over $\mathbb Z$. Under evaluation in
any associative unital ring integer scalars map to central elements,
so the exact same circuit remains valid. No target-ring inverse or
field-extension constant is used.
\end{proof}

\paragraph{Inherited output reductions.}
The degree-four distinct-index polynomial
$S_{4,N}=\sum_{i_1,i_2,i_3,i_4\text{ distinct}}X_{i_1}X_{i_2}X_{i_3}X_{i_4}$
and the row-ordered rectangular polynomial
$\operatorname{perm}_{4,N}=\sum_{i_1,i_2,i_3,i_4\text{ distinct}}
 X_{1,i_1}X_{2,i_2}X_{3,i_3}X_{4,i_4}$
have the same bound by the established integer identities of
Hrube\v{s}~\cite[Proposition~26]{hrubes2026}. These are output transfers
of the same construction, not separate algorithmic contributions.
For completeness put $p_j=\sum_i X_i^j$, $ID^*=ID_N(X,X)$, and
\[
 S_2=p_1^2-p_2,\qquad
 S_3=p_1S_2-p_2p_1-\sum_i X_i p_1 X_i+2p_3.
\]
Then $S_4=p_1S_3-A_0-B_0-C_0$, where
\begin{align*}
 A_0&=p_2S_2-p_3p_1-\sum_iX_i^2p_1X_i+2p_4,\\
 B_0&=\sum_i X_i p_1X_i p_1-ID^*
       -\sum_iX_i p_1X_i^2-p_3p_1+2p_4,\\
 C_0&=\sum_iX_iS_2X_i-\sum_iX_i^2p_1X_i
       -\sum_iX_i p_1X_i^2+2p_4.
\end{align*}
Every sum outside $ID^*$ has linear circuit size, with integer
coefficients. Finally substitute for each $X_j$ the strictly
upper-shift $5\times5$ matrix whose $(r,r+1)$ entry is $X_{r,j}$.
The $(1,5)$ entry is exactly $\operatorname{perm}_{4,N}$.
Matrix simulation costs a constant factor per gate and needs no
division. This concerns four rows; it is not the ordinary growing-order
square permanent.
